% ============================================================
% Fiche de révision — Ch05 Équations du premier degré
% 2nde Bac Pro MAMA — Mathématiques
% Compilable : pdflatex (Overleaf)
% ============================================================
\documentclass[a4paper,10pt]{article}

% --- Encodage et langue ---
\usepackage[utf8]{inputenc}
\usepackage[T1]{fontenc}
\usepackage[french]{babel}

% --- Mise en page compacte (1 page A4) ---
\usepackage[margin=1.2cm,top=1cm,bottom=1cm]{geometry}
\pagestyle{empty}
\setlength{\parindent}{0pt}
\setlength{\parskip}{2pt}

% --- Colonnes ---
\usepackage{multicol}
\setlength{\columnsep}{14pt}

% --- Couleurs et boîtes ---
\usepackage[most]{tcolorbox}
\usepackage{xcolor}

% Palette cohérente avec le site (maths/seconde)
\definecolor{prim}{HTML}{0056b3}
\definecolor{primbg}{HTML}{ebf5ff}
\definecolor{primborder}{HTML}{bee3f8}
\definecolor{greenbox}{HTML}{dcfce7}
\definecolor{greenborder}{HTML}{22c55e}
\definecolor{redbox}{HTML}{fee2e2}
\definecolor{redborder}{HTML}{ef4444}
\definecolor{orangebox}{HTML}{fff7ed}
\definecolor{orangeborder}{HTML}{f97316}
\definecolor{graybg}{HTML}{f8fafc}

% --- Mathématiques ---
\usepackage{amsmath,amssymb}

% --- Divers ---
\usepackage{enumitem}
\setlist{nosep,leftmargin=14pt}

% ============================================================
% Styles tcolorbox
% ============================================================

% En-tête titre
\newtcolorbox{fichetitre}{
  colback=prim,colframe=prim,
  colupper=white,fontupper=\bfseries\large,
  boxrule=0pt,arc=4pt,
  left=6pt,right=6pt,top=4pt,bottom=4pt
}

% Définition
\newtcolorbox{defbox}[1][Définition]{
  colback=primbg,colframe=primborder,
  fonttitle=\bfseries\small,title=#1,
  boxrule=0.8pt,arc=3pt,
  left=5pt,right=5pt,top=3pt,bottom=3pt,
  toptitle=2pt,bottomtitle=2pt
}

% Formule
\newtcolorbox{formulebox}[1][Formule]{
  colback=primbg,colframe=prim,
  fonttitle=\bfseries\small,title=#1,
  boxrule=1pt,arc=3pt,
  left=5pt,right=5pt,top=3pt,bottom=3pt,
  toptitle=2pt,bottomtitle=2pt
}

% Méthode
\newtcolorbox{methbox}[1][Méthode]{
  colback=orangebox,colframe=orangeborder,
  fonttitle=\bfseries\small,title=#1,
  boxrule=0.8pt,arc=3pt,
  left=5pt,right=5pt,top=3pt,bottom=3pt,
  toptitle=2pt,bottomtitle=2pt
}

% Piège (attention)
\newtcolorbox{piegebox}[1][Piège]{
  colback=redbox,colframe=redborder,
  fonttitle=\bfseries\small,title={\color{redborder}$\boldsymbol{\times}$ #1},
  boxrule=0.8pt,arc=3pt,
  left=5pt,right=5pt,top=2pt,bottom=2pt,
  toptitle=2pt,bottomtitle=2pt
}

% Astuce (conseil)
\newtcolorbox{astucebox}[1][Astuce]{
  colback=greenbox,colframe=greenborder,
  fonttitle=\bfseries\small,title={\color{greenborder}$\checkmark$ #1},
  boxrule=0.8pt,arc=3pt,
  left=5pt,right=5pt,top=2pt,bottom=2pt,
  toptitle=2pt,bottomtitle=2pt
}

% Résumé express
\newtcolorbox{resumebox}{
  colback=graybg,colframe=prim,
  boxrule=1.2pt,arc=4pt,
  left=6pt,right=6pt,top=4pt,bottom=4pt
}

% ============================================================
\begin{document}

% --- TITRE ---
\begin{fichetitre}
\centering Fiche de révision — Ch05 — Équations du premier degré\\[1pt]
{\normalsize\normalfont\color{white} Seconde Bac Pro \,$\vert$\, Mathématiques}
\end{fichetitre}

\vspace{4pt}

% ============================================================
% DEUX COLONNES : Définition + Formules
% ============================================================
\begin{multicols}{2}

% --- DÉFINITION ---
\begin{defbox}[Définition]
Une \textbf{équation du premier degré} à une inconnue est une égalité contenant une inconnue $x$ qui n'apparaît qu'au \textbf{degré~1}.

\smallskip
\textbf{Formes courantes :}
\begin{itemize}
  \item $ax = b$
  \item $ax + b = c$
  \item $ax + b = cx + d$
\end{itemize}

\smallskip
{\small $a, b, c, d$ : nombres connus \quad $x$ : inconnue à trouver.}
\end{defbox}

\vspace{4pt}

% --- VOCABULAIRE ---
\begin{defbox}[Vocabulaire]
\begin{itemize}
  \item \textbf{Inconnue} : quantité cherchée ($x$)
  \item \textbf{Solution} : valeur de $x$ qui vérifie l'égalité
  \item \textbf{Membre gauche / droit} : à gauche / droite du signe $=$
\end{itemize}
\end{defbox}

\columnbreak

% --- FORMULES ---
\begin{formulebox}[Formules de résolution]
\[
  ax = b \;\Longrightarrow\; x = \frac{b}{a} \quad (a \neq 0)
\]
\[
  ax + b = c \;\Longrightarrow\; x = \frac{c - b}{a}
\]

\smallskip
\textbf{Règles d'équivalence :}
\begin{itemize}
  \item On peut \textbf{ajouter} ou \textbf{soustraire} le même nombre des deux côtés.
  \item On peut \textbf{multiplier} ou \textbf{diviser} les deux côtés par un nombre \textbf{non nul}.
  \item Un terme qui \textbf{change de côté} change de \textbf{signe}.
\end{itemize}
\end{formulebox}

\end{multicols}

\vspace{2pt}

% ============================================================
% MÉTHODE : 5 étapes
% ============================================================
\begin{methbox}[Méthode — Les 5 étapes de résolution d'un problème]
\begin{enumerate}[label=\textbf{\arabic*.},leftmargin=18pt]
  \item \textbf{Choisir l'inconnue} : nommer $x$, préciser son unité.
  \item \textbf{Traduire l'énoncé} : écrire l'équation à partir des conditions du problème.
  \item \textbf{Résoudre} : isoler $x$ étape par étape (soustraire, puis diviser).
  \item \textbf{Vérifier} : remplacer $x$ par la valeur trouvée dans l'\textbf{équation de départ}.
  \item \textbf{Conclure} : répondre en phrase avec l'unité et interpréter dans le contexte.
\end{enumerate}
\end{methbox}

\vspace{2pt}

% ============================================================
% PIÈGES (3 boîtes rouges)
% ============================================================
\begin{multicols}{3}

\begin{piegebox}[Piège 1 — Mauvaise traduction]
{\small Bien identifier \textbf{ce qui est égal à quoi}. Ne pas confondre somme et produit.

\medskip
\textit{Ex :} « 3 planches et 10\,€ de frais = 55\,€ »

$\rightarrow\; 3x + 10 = 55$

$\not\rightarrow\; 3x \times 10 = 55$}
\end{piegebox}

\columnbreak

\begin{piegebox}[Piège 2 — Signe moins perdu]
{\small Quand on déplace un terme, il \textbf{change de signe}.

\medskip
$5x + 3 = -12$

$5x = -12 \mathbf{-\,3} = -15$

\textbf{et non pas :}

$5x = -12 + 3 = -9$~\rlap{$\;\boldsymbol{\times}$}}
\end{piegebox}

\columnbreak

\begin{piegebox}[Piège 3 — Pas de conclusion]
{\small Ne pas s'arrêter à $x = 4{,}7$.

\medskip
\textit{Ex :} $x = 4{,}7$ planches $\rightarrow$ on ne peut pas couper 4,7~planches !

$\rightarrow$ \textbf{Arrondir} (4 ou 5) et \textbf{interpréter} dans le contexte.}
\end{piegebox}

\end{multicols}

\vspace{2pt}

% ============================================================
% ASTUCES (2 boîtes vertes)
% ============================================================
\begin{multicols}{2}

\begin{astucebox}[Astuce 1 — Vérification]
Toujours vérifier dans l'\textbf{équation de départ}, jamais dans une étape intermédiaire (qui pourrait contenir une erreur).

\smallskip
{\small\textit{Ex :} si $3x + 5 = 20$ donne $x = 5$, vérifier $3 \times 5 + 5 = 20$~\checkmark}
\end{astucebox}

\columnbreak

\begin{astucebox}[Astuce 2 — Résolution graphique]
La solution de $ax + b = c$ est l'\textbf{abscisse du point d'intersection} des droites $y = ax + b$ et $y = c$.

\smallskip
{\small Utile pour vérifier avec la calculatrice graphique ou GeoGebra.}
\end{astucebox}

\end{multicols}

\vspace{4pt}

% ============================================================
% RÉSUMÉ EXPRESS
% ============================================================
\begin{resumebox}
\centering
\textbf{Résumé express}\\[4pt]
\fbox{\strut\textbf{Poser $x$}}
$\;\longrightarrow\;$
\fbox{\strut\textbf{Écrire l'équation}}
$\;\longrightarrow\;$
\fbox{\strut\textbf{Regrouper les termes}}
$\;\longrightarrow\;$
\fbox{\strut\textbf{Diviser par le coefficient}}
$\;\longrightarrow\;$
\fbox{\strut\textbf{Vérifier \& Conclure}}
\end{resumebox}

\end{document}
